\documentclass[11pt]{article}
\usepackage[margin=1in]{geometry}
\usepackage{amsmath,amssymb}
\usepackage{graphicx}
\usepackage{xcolor}
\usepackage{tcolorbox}
\usepackage{enumitem}
\usepackage{hyperref}
\usepackage{titlesec}
\usepackage{fancyhdr}
\usepackage{booktabs}
\usepackage{parskip}

\definecolor{navy}{HTML}{1E2761}
\definecolor{ice}{HTML}{CADCFC}
\definecolor{red}{HTML}{E63946}
\definecolor{lightgrey}{HTML}{F4F5F9}

\tcbuselibrary{skins,breakable}

\titleformat{\section}{\Large\bfseries\color{navy}}{\thesection}{0.5em}{}
\titleformat{\subsection}{\large\bfseries\color{navy}}{\thesubsection}{0.5em}{}

\pagestyle{fancy}
\fancyhf{}
\fancyhead[L]{\small\color{gray} Study Guide: Data Balancing for Hotspot Detection}
\fancyhead[R]{\small\color{gray} Group 11}
\fancyfoot[C]{\small\color{gray} \thepage}
\renewcommand{\headrulewidth}{0.4pt}

\newtcolorbox{keyidea}[1][]{
  colback=lightgrey, colframe=navy, boxrule=0pt, leftrule=3pt,
  arc=2pt, outer arc=2pt, fonttitle=\bfseries, title={#1},
  coltitle=white, breakable
}
\newtcolorbox{analogy}[1][Analogy]{
  colback=ice!30, colframe=navy!60, boxrule=0pt, leftrule=3pt,
  arc=2pt, outer arc=2pt, fonttitle=\bfseries\itshape, title={#1},
  coltitle=white, breakable
}
\newtcolorbox{takeaway}[1][Takeaway]{
  colback=red!6, colframe=red, boxrule=0pt, leftrule=3pt,
  arc=2pt, outer arc=2pt, fonttitle=\bfseries, title={#1},
  coltitle=white, breakable
}

\title{\vspace{-1.5cm}\textbf{\color{navy}\Huge Study Guide}\\[0.3cm]
\Large Data Balancing and Augmentation Strategies\\ for Lithography Hotspot Detection\\[0.3cm]
\normalsize\itshape A beginner's walkthrough of the AI/ML concepts and the project itself, from zero prior knowledge}
\author{Written for: Group 11 (Priyanuj Boruah, Kartik Aggarwal)\\ \small BEVD402L -- AI and Machine Learning for IC}
\date{}

\begin{document}
\maketitle
\thispagestyle{empty}

\begin{keyidea}[How to use this guide]
This is not the assignment report -- it is a plain-language explainer,
written assuming you know \emph{nothing} about AI/ML going in. Read it
top to bottom once, then use it as a reference when someone asks
``wait, what does that actually mean?'' about any part of the project.
Every technical term is introduced with a real-world analogy before
it's used seriously.
\end{keyidea}

\newpage
\tableofcontents
\newpage

\section{Part 1: AI and Machine Learning, From Zero}

\subsection{What is ``Machine Learning'' really?}
Ordinary computer programs work from rules a human writes: \emph{``if
the temperature is above 30, turn on the fan.''} Machine learning (ML)
flips this around. Instead of a human writing the rule, you show the
computer thousands of examples, and it \emph{figures out the rule
itself}.

\begin{analogy}
Imagine teaching a child to tell cats from dogs. You don't hand them a
rulebook (``if it has pointy ears and whiskers and is under 40cm
long...''). You just show them hundreds of photos, say ``cat'' or
``dog'' each time, and eventually they start guessing correctly on
photos they've never seen. That's exactly what machine learning does
-- except the ``child'' is a mathematical model, and the ``photos''
are training data.
\end{analogy}

\subsection{The two phases: Training and Testing}
Every ML project has two distinct phases, and mixing them up is the
single most common beginner mistake:

\begin{itemize}[leftmargin=1.2em]
    \item \textbf{Training} -- the model looks at examples \emph{with
    the correct answer attached} (called \emph{labels}) and adjusts
    itself to get better at predicting that answer.
    \item \textbf{Testing (or evaluation)} -- the model is shown
    \emph{new} examples it has never seen before, \emph{without} the
    answer, and we check how many it gets right.
\end{itemize}

\begin{keyidea}[Why the split matters]
If you test the model on the same examples it trained on, it can
simply ``memorize'' the answers instead of actually learning the
underlying pattern -- like a student who memorizes yesterday's exam
questions verbatim instead of understanding the topic. That's called
\textbf{overfitting}, and it's why every serious ML project keeps
training data and test data strictly separate. In this project, this
distinction is critical: we only ever apply our balancing/augmentation
tricks to the \emph{training} data; the test data is always left
exactly as it was found.
\end{keyidea}

\subsection{What is a Neural Network?}
A neural network is one particular kind of ML model, loosely inspired
by how brain cells (neurons) connect to each other. Don't worry about
the biology -- think of it as a big mathematical function made of
layers of simple building blocks, where:

\begin{itemize}[leftmargin=1.2em]
    \item Each ``neuron'' takes some numbers in, multiplies them by
    learned weights, adds them up, and passes the result through a
    simple non-linear switch (called an \textbf{activation
    function}).
    \item Layers of neurons are stacked, so the output of one layer
    feeds into the next.
    \item \textbf{Training} means slowly adjusting all those weights
    (there can be thousands to billions of them) so that the
    network's final output matches the correct label as often as
    possible.
\end{itemize}

\begin{analogy}
Think of a factory assembly line with many stations. Raw material
(the image) enters at one end. Each station does a small
transformation -- highlighting edges, detecting shapes, combining
patterns -- and passes its output to the next station. By the final
station, the ``product'' is a single decision: \emph{hotspot} or
\emph{not a hotspot}. Training is the process of tuning every station
on the line so the final product is usually correct.
\end{analogy}

\subsection{What is a CNN (Convolutional Neural Network)?}
A CNN is a neural network specialized for images. The key building
block is a \textbf{convolution}: a small sliding window (e.g.\ 3
pixels by 3 pixels) that scans across the image looking for a specific
small pattern -- an edge, a corner, a line of a certain thickness.
Many such windows (called \textbf{filters}) run in parallel, each
learning to detect a different small pattern.

\begin{itemize}[leftmargin=1.2em]
    \item \textbf{Convolution layer:} slides small filters over the
    image to detect local patterns (edges, corners, textures).
    \item \textbf{Pooling layer:} shrinks the image down (e.g.\ takes
    the strongest signal in every 2$\times$2 block), which both
    reduces computation and makes the network care less about exactly
    \emph{where} a pattern appeared.
    \item \textbf{Fully-connected layer:} at the very end, takes all
    the detected patterns and combines them into a final decision.
\end{itemize}

Stacking several convolution+pooling stages lets early layers detect
simple things (edges) and later layers combine those into complex
things (a ``pinched wire'' shape that indicates a hotspot).

\begin{keyidea}[Why CNNs, and not something simpler, for this project]
A layout image is just a grid of pixels. A CNN is good at this task
because a hotspot is defined by a \emph{local geometric pattern} (two
lines getting too close together, a corner that's too sharp) --
exactly the kind of small, spatial pattern that convolution filters
are built to detect, wherever it appears in the image.
\end{keyidea}

\subsection{Binary Classification}
Our task has exactly two possible answers: \emph{Hotspot (HS)} or
\emph{Non-Hotspot (NHS)}. This is called \textbf{binary
classification} (``binary'' = two options). The network's very last
layer outputs a single number between 0 and 1 (think of it as ``how
confident am I that this is a hotspot, from 0\% to 100\%''), and we
typically call it a hotspot if that number is above 0.5.

\subsection{Precision, Recall, and Why ``Accuracy'' Can Lie To You}
This is the single most important idea for understanding this entire
project, so we'll spend real time on it.

\textbf{Accuracy} is the simplest possible score: what fraction of all
predictions were correct? It sounds reasonable, but it has a fatal
flaw when one class is much rarer than the other -- which is exactly
our situation (see Part 3).

\begin{analogy}
Imagine a medical test for a rare disease that only 1 in 1,000 people
actually have. A ``test'' that always says ``you're healthy,'' no
matter what, will be correct 999 times out of 1,000 -- 99.9\%
accuracy! But it is completely useless: it will never, ever catch the
one person who is actually sick. High accuracy, zero usefulness.
\end{analogy}

This is precisely the trap in lithography hotspot detection: real
hotspots are rare (Part 3), so a model that just says ``not a
hotspot'' every single time can score above 99\% raw accuracy on some
of our benchmarks while catching \emph{zero} real hotspots. That is
why this whole project uses different scores instead:

\begin{itemize}[leftmargin=1.2em]
    \item \textbf{Recall (a.k.a.\ Sensitivity):} of all the
    \emph{actual} hotspots, what fraction did we correctly catch?
    This is the ``did we miss the disease'' score.
    \item \textbf{Specificity:} of all the actual \emph{non}-hotspots,
    what fraction did we correctly leave alone? This is the ``did we
    cry wolf'' score.
    \item \textbf{Precision:} of everything we \emph{flagged} as a
    hotspot, what fraction actually were? High precision means few
    false alarms.
    \item \textbf{Balanced Accuracy:} the average of Recall and
    Specificity. This is the project's main scorecard, precisely
    because it can't be gamed by just always predicting the majority
    class -- to score well, a model must do well on \emph{both}
    classes.
    \item \textbf{F1-score:} a combined score balancing Precision and
    Recall, useful when you care about both false alarms and missed
    detections together.
\end{itemize}

\begin{keyidea}[The Confusion Matrix]
All of the above scores come from just four numbers, usually laid out
in a $2\times2$ grid called a \textbf{confusion matrix}:
\begin{center}
\begin{tabular}{ll|c|c}
 & & \multicolumn{2}{c}{\textbf{Model Predicted}} \\
 & & Hotspot & Non-Hotspot \\ \hline
\textbf{Actually} & Hotspot & True Positive (TP) & False Negative (FN) \\
\textbf{Was} & Non-Hotspot & False Positive (FP) & True Negative (TN) \\
\end{tabular}
\end{center}
FN (a real hotspot we missed) is usually the most dangerous error in
this domain -- it slips through to manufacturing. FP (a false alarm)
just costs some wasted verification time.
\end{keyidea}

\section{Part 2: The Real-World Problem -- Lithography Hotspots}

\subsection{How is a computer chip made?}
A modern chip has billions of microscopic wires and transistors
packed onto a few square centimeters of silicon. The pattern of these
wires is ``printed'' onto the silicon using a process called
\textbf{photo-lithography} -- shining light through a stencil (called
a \textbf{mask}) coated with the desired circuit pattern, similar in
spirit to how a photographic negative projects an image, but at a
scale of billionths of a metre.

\subsection{What is a ``hotspot''?}
As chip features get smaller and smaller, they start to approach the
same size as the \emph{wavelength of light} being used to print them.
When a wire (or the gap between two wires) is on the edge of what the
light can reliably resolve, small pattern variations can cause a real
manufacturing defect:

\begin{itemize}[leftmargin=1.2em]
    \item \textbf{Pinching (open circuit):} two ends of a wire that
    were supposed to connect end up separated -- like a bridge that
    doesn't quite meet in the middle.
    \item \textbf{Bridging (short circuit):} two separate wires that
    were supposed to stay apart end up merging together -- like two
    lanes of a highway accidentally merging.
\end{itemize}

A \textbf{hotspot} is a small region of the chip layout that is at
high risk of one of these failures once actually printed. Figure~1 of
the reference paper shows real examples: a design's intended pattern
(``Original'') next to what actually got printed (``Printed''), with
the defect circled.

\begin{keyidea}[Why we bother detecting these early]
Finding a hotspot \emph{after} manufacturing means the chip is
already defective -- expensive, wasted silicon. Finding it during
\emph{design}, before sending the layout to the fab, means engineers
can just tweak that one small region and re-check, which is far
cheaper. The gold-standard way to check is full optical simulation of
how light will behave -- but that is slow and computationally very
expensive to run across an entire chip. This project (like the wider
field it belongs to) is about using a fast, approximate, ML-based
screening step to flag \emph{likely} hotspots quickly, so full
simulation only needs to double-check the suspicious regions.
\end{keyidea}

\subsection{The ICCAD-12 Dataset}
To let researchers compare their methods fairly, a standard
benchmark dataset was released in 2012, called \textbf{ICCAD-12}. It
contains five separate collections (``Benchmark 1'' through
``Benchmark 5''), each a set of small image clips cut from real chip
layouts, where every clip is labeled either Hotspot (HS) or
Non-Hotspot (NHS). This is the exact dataset our whole project uses,
and it is what makes results from different student groups (and
different research papers) directly comparable.

\section{Part 3: Our Specific Problem -- Class Imbalance}

\subsection{What does ``imbalanced'' mean here?}
In the real world, the overwhelming majority of any chip design is
completely fine -- hotspots are, by definition, rare exceptions. This
shows up directly in the dataset:

\begin{center}
\begin{tabular}{lccc}
\toprule
Benchmark & Hotspot images & Non-Hotspot images & Ratio \\
\midrule
B1 & 226  & 4{,}679  & 1 hotspot per 21 images \\
B2 & 498  & 41{,}298 & 1 hotspot per 83 images \\
B3 & 1{,}808 & 46{,}333 & 1 hotspot per 26 images \\
B4 & 177  & 31{,}890 & 1 hotspot per 180 images \\
B5 & 41   & 19{,}327 & 1 hotspot per \textbf{471} images \\
\bottomrule
\end{tabular}
\end{center}

On Benchmark 5, for every single real hotspot, there are 471
non-hotspot images. This is what we mean by \textbf{severe class
imbalance}, and it is the central obstacle this whole project is
about.

\begin{analogy}
Imagine trying to teach someone to recognize a rare bird species by
showing them 471 pictures of pigeons for every 1 picture of the rare
bird. If they just learned to say ``pigeon'' every single time, they'd
be right 99.8\% of the time -- and would never learn to recognize the
rare bird at all. A neural network trained naively on this kind of
data has exactly the same temptation: predicting the majority class
every time is a very easy way to rack up high raw accuracy.
\end{analogy}

\subsection{Why does this matter for training, specifically?}
Neural networks learn by seeing many examples and gradually reducing
their mistakes. If 99.8\% of what they see is one class, the easiest
way to reduce mistakes is to just lean heavily towards predicting that
class -- the rare class barely gets enough ``votes'' during training
to matter. Our project's whole job is to fix this imbalance in the
\emph{training data} so the model actually has to learn what a real
hotspot looks like.

\section{Part 4: The Six Strategies We Compared}
This is Group 11's specific research angle: \emph{comparing different
ways of fixing the training data's imbalance}, to see which works
best. Here is each strategy explained with a plain-language analogy.

\begin{keyidea}[0. Baseline -- Do Nothing]
Train on the data exactly as it comes, imbalance and all. This is our
reference point -- every other strategy is judged by how much (or how
little) it improves on this.
\end{keyidea}

\begin{keyidea}[1. Random Undersampling (RUS)]
\textbf{Idea:} Throw away most of the majority (non-hotspot) examples
until the two classes are more equal in number.
\begin{analogy}
If you have 4,700 pigeon photos and 226 rare-bird photos, just
randomly delete pigeon photos until you have roughly 226 of each. Now
the training set is balanced -- but you've thrown away a lot of
pigeon photos you could have learned from, and you haven't added a
single new rare-bird photo.
\end{analogy}
\textbf{Risk:} If the rare class had very few examples to begin with,
you end up with a tiny total training set, and the model may not have
enough information to learn from at all.
\end{keyidea}

\begin{keyidea}[2. Random Oversampling (ROS)]
\textbf{Idea:} Instead of deleting majority examples, duplicate
minority (hotspot) examples until the classes are more equal.
\begin{analogy}
Instead of deleting pigeon photos, make 20 identical copies of each
rare-bird photo you have, so the count matches. You keep everything,
but the rare-bird photos are now exact repeats.
\end{analogy}
\textbf{Risk:} The model might start ``memorizing'' those exact
duplicated images rather than learning the general pattern --
overfitting to a small handful of examples.
\end{keyidea}

\begin{keyidea}[3. Oversampling + Augmentation]
\textbf{Idea:} Same as ROS, but instead of making \emph{identical}
copies, apply small, label-preserving changes to each duplicate --
flipping it left-right, flipping it top-bottom, or rotating it 90/180/
270 degrees.
\begin{analogy}
Instead of 20 identical photocopies of the same rare bird photo, you
make 20 copies but mirror some, flip some upside down, and rotate
some. It's still recognizably the same bird, but the model sees more
\emph{variety}, making it a little harder to just memorize pixels.
\end{analogy}
\textbf{Why flips/rotations specifically?} A hotspot is still a
hotspot if you flip or rotate the whole image -- it's the same
physical pattern, just viewed from a different orientation, so this
doesn't create fake/incorrect examples.
\end{keyidea}

\begin{keyidea}[4. Balanced Mini-Batch Sampling]
\textbf{Idea:} Don't touch the dataset at all. Instead, change how
training examples are \emph{picked} in each small batch during
training, so that every batch shown to the model has roughly equal
hotspot and non-hotspot examples.
\begin{analogy}
Instead of editing your photo collection, you change how you
\emph{study} it: every study session, you deliberately pull an equal
number of pigeon and rare-bird photos from your (untouched, full)
collection. No photo was deleted or copied -- you just changed your
study habits.
\end{analogy}
\textbf{Why this often works well:} it gets the benefit of balanced
exposure during training without physically discarding data (like
RUS) or risking memorized duplicates (like ROS).
\end{keyidea}

\begin{keyidea}[5. Combined Strategy]
\textbf{Idea:} Do a bit of both -- moderately reduce the majority
class \emph{and} moderately increase the (augmented) minority class,
meeting somewhere in the middle, rather than going to either extreme.
\end{keyidea}

\section{Part 5: What We Actually Built}

\begin{enumerate}[leftmargin=1.3em]
    \item \textbf{Got the real data.} Downloaded the actual ICCAD-12
    dataset (all five benchmarks, tens of thousands of real layout
    images) and confirmed the hotspot/non-hotspot counts matched the
    published reference paper.
    \item \textbf{Built one shared ``brain''.} Implemented a single
    lightweight CNN (about 8,600 tunable numbers -- tiny by modern AI
    standards, deliberately so it trains fast) that every strategy
    uses identically. This is important: if every strategy uses the
    exact same network, any performance difference we see must come
    from the \emph{training data} itself, not from a better or worse
    network design. That is what makes this a fair, controlled
    experiment.
    \item \textbf{Implemented all six strategies} as code, each one
    taking the raw imbalanced training images in and handing back a
    (possibly resampled/augmented) training set, never touching the
    test images.
    \item \textbf{Trained and tested 105 times.} 5 benchmarks
    $\times$ 6 strategies = 30 main experiments, plus 75 more
    experiments where we varied \emph{how strongly} to rebalance the
    data (to see if going halfway is better than going all the way).
    \item \textbf{Measured everything that matters.} For every one of
    those 105 runs, we recorded the full confusion matrix, balanced
    accuracy, precision, recall, specificity, and F1-score -- never
    just raw accuracy, for the reasons explained in Part 1.
    \item \textbf{Turned the numbers into pictures and a paper.}
    Generated charts and tables, and wrote them up as a proper
    IEEE-style research report.
\end{enumerate}

\section{Part 6: What We Found -- Explained Simply}

\subsection{Headline result}
Averaged across all five benchmarks:

\begin{center}
\begin{tabular}{lc}
\toprule
Strategy & Average Balanced Accuracy \\
\midrule
Balanced Mini-Batch Sampling & \textbf{0.945 (best)} \\
Combined & 0.926 \\
Random Oversampling & 0.923 \\
Oversampling + Augmentation & 0.887 \\
Baseline (do nothing) & 0.844 \\
Random Undersampling & \textbf{0.779 (worst)} \\
\bottomrule
\end{tabular}
\end{center}

In plain terms: \textbf{just re-sampling which examples get shown to
the model in each training batch (Strategy \#4) worked better, on
average, than any other trick} -- including fancier ones that alter
the dataset itself. And somewhat surprisingly, \textbf{the strategy
that most directly ``fixes'' the imbalance by deleting data
(undersampling) was actually the single worst performer.}

\subsection{The most dramatic finding: Benchmark 5}
Benchmark 5 is the extreme case -- only 26 real hotspot training
images, against 2,716 non-hotspot ones. Here's what happened:

\begin{itemize}[leftmargin=1.2em]
    \item \textbf{Random Undersampling:} caught \textbf{0} of 41 real
    hotspots in the test set. It predicted ``not a hotspot'' every
    single time, no matter what we fed it -- even though we had
    artificially balanced its training set to 50/50! We even tried
    every possible balance level (from mild to full) and it stayed
    stuck at 0 hotspots caught, every time.
    \item \textbf{Balanced Mini-Batch Sampling:} caught \textbf{40 of
    41} real hotspots (97.6\% recall).
\end{itemize}

\begin{takeaway}
Why did undersampling fail so badly? Because it only ever had 26
\emph{original} hotspot images to work with -- deleting non-hotspot
images doesn't create a single new piece of information about what a
hotspot looks like. If you only ever show a child 26 photos of a rare
bird, no amount of throwing away pigeon photos will suddenly teach
them what that bird looks like. The lesson: \textbf{when data is
critically scarce, the fix has to add information (or at least
preserve all of it), not just remove the competing class.}
\end{takeaway}

\subsection{A genuinely surprising finding: augmentation can backfire}
We expected ``oversampling + augmentation'' (Strategy \#3) to always
beat plain oversampling (Strategy \#2), since flipping/rotating adds
variety. That was true on Benchmarks 2 and 3 (augmentation clearly
helped). But on Benchmark 5 -- the one with only 26 original hotspot
images -- augmentation actually made things \emph{worse} than plain
duplication (balanced accuracy 0.707 vs.\ 0.914).

\begin{takeaway}
When you have very few original examples, almost every hotspot image
the network sees during training ends up being a flipped/rotated copy
of one of those few originals. The model can end up learning quirks
of the flipping/rotating process itself, rather than the real
geometric pattern of a hotspot -- and that hurts it on the real
(unflipped) test images. \textbf{The lesson: a technique that helps in
one setting can hurt in another -- there's no single ``best''
strategy that works everywhere, which is exactly the nuanced answer
our research question was looking for.}
\end{takeaway}

\subsection{What the confusion matrix numbers mean, concretely}
On Benchmark 2, the do-nothing baseline caught only 197 of 498 real
hotspots (missed 301!) while being extremely cautious about false
alarms. Switching to the Combined strategy caught 495 of 498 real
hotspots, at the cost of a modest number of extra false alarms (415
non-hotspots incorrectly flagged, out of 41,713 total). In a
real manufacturing context, missing 301 real defects is a far more
expensive mistake than double-checking 415 patterns that turned out
to be fine -- which is exactly the kind of trade-off Balanced Accuracy
(and the confusion matrix) is designed to make visible.

\section{Glossary (Quick Reference)}
\begin{description}[leftmargin=1.6em, style=nextline]
\item[Accuracy] Fraction of all predictions that were correct. Misleading when classes are imbalanced.
\item[Balanced Accuracy] Average of Recall and Specificity; the fair scorecard used throughout this project.
\item[Baseline] Training with no balancing intervention -- the reference point.
\item[Class imbalance] When one label (here, non-hotspot) vastly outnumbers the other (hotspot) in the data.
\item[CNN (Convolutional Neural Network)] A neural network specialized for images, using sliding filters to detect local visual patterns.
\item[Confusion matrix] The 2$\times$2 table of TP/FP/FN/TN counts that every other metric is computed from.
\item[Data augmentation] Creating variations of existing examples (e.g.\ flips, rotations) to add diversity without collecting new data.
\item[F1-score] Combined score balancing Precision and Recall.
\item[Hotspot (HS)] A layout region at high risk of a real manufacturing defect once printed.
\item[Non-Hotspot (NHS)] A layout region that is expected to print correctly.
\item[Oversampling] Duplicating minority-class examples to reduce imbalance.
\item[Overfitting] When a model memorizes training examples instead of learning general patterns, and performs worse on new data.
\item[Precision] Of everything flagged positive, what fraction was correct.
\item[Recall (Sensitivity)] Of all real positives, what fraction was correctly caught.
\item[Specificity] Of all real negatives, what fraction was correctly left alone.
\item[Test set] Data held back and never used for training, used only to check final performance.
\item[Training set] Data (with correct labels) shown to the model so it can learn.
\item[Undersampling] Discarding majority-class examples to reduce imbalance.
\end{description}

\end{document}
